\section*{Chapitre \ref{ch:iv2:options-and-bond-arithmetic-at-speed} --- Options et arithmétique obligataire à toute vitesse}
\markboth{\protect\omTocMark{}{Réponses~: Options et arithmétique obligataire à toute vitesse}}{\protect\omTocMark{}{Réponses~: Options et arithmétique obligataire à toute vitesse}}

\begin{solution}{iq:iv2:options-and-bond-arithmetic-at-speed:1}
Environ 8~\$ et environ 4.24~\$. La règle dit qu'une option d'achat à la monnaie vaut environ $0.4\,S\sigma\sqrt{T}$. Les valeurs exactes de Black-Scholes sont 7.97~\$ et 4.22~\$, si bien que la règle est précise à mieux que 1~\% pour ces volatilités.

Ici, 0.4 désigne $1/\sqrt{2\pi}$. Lorsque les taux sont nuls et que le strike est égal au spot, le prix de Black-Scholes\index{modèle de Black-Scholes} vaut $S(2\Phi(\sigma\sqrt T/2) - 1)$, et au premier ordre c'est $S \cdot 2\varphi(0) \cdot \sigma\sqrt T/2$. Donc $0.4 \times 100 \times 0.2 = 8$, et $0.4 \times 50 \times 0.3 \times 0.707 \approx 4.24$.

\iqwrong{Utiliser $\sqrt{T}$ avec $T$ exprimé en mois, ou oublier la racine carrée. Le temps met la volatilité à l'échelle par sa racine carrée, si bien que six mois représentent 71~\% de la volatilité annuelle, et non la moitié.}
\end{solution}

\begin{solution}{iq:iv2:options-and-bond-arithmetic-at-speed:2}
D'après $0.8\,S\sigma\sqrt T = 0.8 \times 200 \times 0.4 \times 0.5$, le straddle coûte environ 32~\$. La valeur exacte est d'environ 31.86~\$. Il atteint son seuil de rentabilité à environ 32~\$ dans un sens comme dans l'autre, soit 16~\%. Un mouvement quotidien typique vaut environ $40~\%/16 = 2.5~\%$.

La valeur du straddle vaut le double de celle de l'option d'achat. C'est aussi le mouvement absolu moyen attendu, puisque $E|Z| = \sqrt{2/\pi} \approx 0.8$ écart-type. Ainsi, le seuil de rentabilité d'un straddle, à 0.8 écart-type, est atteint un peu plus de 40~\% du temps.

Pour un mouvement quotidien, on divise la volatilité annuelle par 16, car $\sqrt{252} \approx 15.9$.

\iqwrong{Diviser par 252 pour obtenir un mouvement quotidien. La variance s'additionne d'un jour à l'autre, si bien que la volatilité croît comme la racine carrée du temps.}
\end{solution}

\begin{solution}{iq:iv2:options-and-bond-arithmetic-at-speed:3}
Le delta de l'option d'achat vaut environ 0.54, et celui de l'option de vente environ $-0.46$. Les valeurs exactes sont $\Phi(0.1) \approx 0.5398$ et $-0.4602$.

Avec des taux nuls et un strike égal au spot, $d_1 = \sigma\sqrt T/2 = 0.1$, et $\Phi(x) \approx 0.5 + 0.4x$ pour $x$ petit. Donc le delta de l'option d'achat vaut $0.5 + 0.04$. Celui de l'option de vente est celui de l'option d'achat diminué de 1, par la parité put-call.

Il est supérieur à un demi parce que le delta est la probabilité de finir dans la monnaie sous une mesure dans laquelle l'action sert d'unité de compte, ce qui donne un poids supplémentaire aux prix élevés. La médiane d'une action lognormale est inférieure à sa moyenne, et cela ramène la probabilité ordinaire en dessous de la moitié~: $\Phi(-0.1) \approx 0.46$.

\iqwrong{«~Une option à la monnaie a un delta exactement égal à 0.5.~» Cela n'est vrai que si la volatilité multipliée par la racine du temps est faible, et que les taux ne font pas remonter le forward au-dessus du strike.}
\end{solution}

\begin{solution}{iq:iv2:options-and-bond-arithmetic-at-speed:4}
Environ 0.40~\$ pour l'option d'achat à un an, environ 0.20~\$ pour l'option d'achat à trois mois, et environ 2~\$ pour l'option d'achat à un an sur l'action à 500~\$.

L'option d'achat vaut $0.4\,S\sigma\sqrt T$, donc le vega\index{vega} vaut $0.4\,S\sqrt T$ par 1.00 de volatilité, et $0.004\,S\sqrt T$ par point de volatilité. Ainsi, à 100~\$ et pour un an, il vaut 0.40, et pour un trimestre, la moitié de cette valeur. Il ne dépend pas de la volatilité elle-même, car à la monnaie le prix est presque linéaire en fonction de la volatilité. La valeur exacte pour une volatilité de 20~\% est d'environ 0.397~\$.

\iqwrong{Penser que les options vendues ont plus de vega parce qu'elles sont «~plus sensibles~». Elles ont plus de gamma. Le vega croît avec la racine carrée du temps.}
\end{solution}

\begin{solution}{iq:iv2:options-and-bond-arithmetic-at-speed:5}
L'option de vente coûte environ 7.12~\$. Rapidement, $12 - 5 = 7$, auquel il faut ajouter une petite correction d'actualisation.

Par la parité put-call\index{parité put-call}, $C - P = S - Ke^{-rT} = 100 - 100e^{-0.05} \approx 4.88$. Donc $P = 12 - 4.88 = 7.12$.

L'estimation rapide utilise le taux multiplié par le strike, soit 5.00. Elle surestime d'environ 0.12, car l'actualisation à 5~\% sur un an retranche un peu moins de 5~\%.

\iqwrong{Dire que l'option de vente et l'option d'achat de même strike doivent coûter autant. Cela n'est vrai que lorsque le strike égale le forward, soit ici $100e^{0.05} \approx 105.13$.}
\end{solution}

\begin{solution}{iq:iv2:options-and-bond-arithmetic-at-speed:6}
La duration seule donne une perte de 487\,500~\$. Avec la convexité, la perte est d'environ 480\,750~\$.

La variation en pourcentage vaut $-D\,\Delta y + \tfrac12 C\,\Delta y^2$. Donc $-6.5 \times 0.30\% = -1.95\%$, soit 487\,500~\$ sur 25 millions de dollars. Le terme de convexité\index{convexité} vaut $\tfrac12 \times 60 \times 0.003^2 = 0.027\%$, soit 6\,750~\$, et il réduit la perte.

\iqwrong{Retrancher la convexité de la variation de prix. La convexité est positive pour une obligation ordinaire, quel que soit le sens de variation des rendements, et elle joue en faveur du détenteur.}
\end{solution}

\begin{solution}{iq:iv2:options-and-bond-arithmetic-at-speed:7}
Le DV01\index{DV01} vaut 8\,200~\$. Il faut vendre environ 126 contrats.

Le DV01 est égal à la valeur multipliée par la duration modifiée et par un point de base~: $10\,000\,000 \times 8.2 \times 0.0001 = 8\,200$. Puis $8\,200/65 \approx 126.2$.

Cette couverture ne protège que contre un mouvement parallèle. Le DV01 d'un contrat à terme est celui de l'obligation qu'il suit, si bien que, si cette obligation se situe à un autre point de la courbe, un risque de courbe subsiste.

\iqwrong{Utiliser la valeur nominale de l'obligation lorsque le prix s'écarte nettement du pair. Le DV01 évolue proportionnellement à la valeur de marché.}
\end{solution}

\begin{solution}{iq:iv2:options-and-bond-arithmetic-at-speed:8}
Le prix vaut environ 67.6. La duration de Macaulay vaut 10, et la duration modifiée environ 9.62. Une hausse de 10 points de base fait baisser le prix d'environ 0.65, pour atteindre environ 66.9.

À la main, $e^{-0.4} \approx 0.67$, avec $1.04$ proche de $e^{0.04}$. La valeur exacte est $100/1.04^{10} = 67.56$. Un zéro-coupon ne comporte qu'un flux, si bien que sa duration de Macaulay est égale à son échéance. La duration modifiée se calcule en divisant par $1 + y$. Puis $9.62 \times 0.1~\% \times 67.56 \approx 0.65$.

\iqwrong{Estimer la duration d'un zéro-coupon à 10 ans à environ 8, comme pour une obligation à coupons. Il n'y a pas de coupons pour ramener la durée de vie moyenne vers le bas.}
\end{solution}

\begin{solution}{iq:iv2:options-and-bond-arithmetic-at-speed:9}
Réponse rapide~: 5~\%. La réponse exacte est d'environ 5.01~\%.

Placer pendant deux ans à 4~\% doit rapporter autant que placer un an à 3~\%, puis un an au taux forward~: $1.04^2 = 1.03(1 + f)$, donc $f = 1.0816/1.03 - 1 \approx 5.01~\%$.

La réponse rapide fonctionne parce que les logarithmes s'additionnent~: le taux à deux ans est à peu près la moyenne du taux à un an et du forward, donc $f \approx 2 \times 4~\% - 3~\%$.

\iqwrong{Faire la moyenne de 3~\% et de 4~\%. Le taux à deux ans est déjà cette moyenne, et c'est le forward qui la produit.}
\end{solution}

\begin{solution}{iq:iv2:options-and-bond-arithmetic-at-speed:10}
Rapidement, $1.10 \times 1.02 = 1.122$. Exactement, $1.10 \times 1.05/1.03 \approx 1.1214$. À 1.13, le forward est trop élevé, il faut donc vendre des euros à terme et les acheter aujourd'hui avec des dollars empruntés.

La devise à taux élevé se traite à un forward en décote, ce qui signifie ici davantage de dollars par euro. On réalise le verrouillage par euro livré dans un an~: on emprunte des dollars aujourd'hui, on achète $1/1.03$ euro à 1.10 pour environ 1.068~\$, et on place ces sommes sur un dépôt. Dans un an, on détient 1 euro, que l'on livre à 1.13. On doit $1.068 \times 1.05 \approx \$1.1214$. Le profit est d'environ 0.86 cent de dollar par euro.

C'est la parité des taux d'intérêt couverte\index{parité des taux d'intérêt couverte}.

\iqwrong{Penser que la devise à taux élevé doit s'apprécier à terme. Le forward compense l'écart de taux, si bien que le dollar, qui a le taux le plus élevé, se négocie à terme plus faiblement face à l'euro.}
\end{solution}

\begin{solution}{iq:iv2:options-and-bond-arithmetic-at-speed:11}
Acheter le panier, créer 50\,000 actions de l'ETF et les vendre au cours acheteur. On gagne 0.09~\$ par action, soit 4\,500~\$ par unité de création.

Vendre à 50.02~\$ contre un coût total de $\$49.90 + \$0.03 = \$49.93$ laisse 0.09~\$. L'ETF est surcoté, donc on le vend et on le crée. Dans le sens inverse, acheter l'ETF à 50.04~\$ et le racheter contre un panier valant 49.90~\$ fait perdre de l'argent.

\iqwrong{Comparer le panier au mid de l'ETF ou à son cours vendeur. On vend au cours acheteur, et les frais se retranchent d'abord.}
\iqfollow{À quel risque êtes-vous exposé~?}{Les cours peuvent bouger pendant que l'on négocie, et le panier n'est pas forcément entièrement négociable à la valeur cotée. On achète donc le panier et on vend l'ETF au même moment, et l'on crée ensuite.}
\end{solution}

\begin{solution}{iq:iv2:options-and-bond-arithmetic-at-speed:12}
La volatilité implicite est d'environ 30~\%. Le mouvement d'un écart-type jusqu'à l'échéance est d'environ 15~\$.

Le straddle coûte $0.8\,S\sigma\sqrt T$, donc $\sigma = 12/(0.8 \times 100 \times 0.5) = 30~\%$. Un écart-type sur trois mois vaut $100 \times 0.3 \times 0.5 = 15$, et $12/0.8 = 15$ donne le même résultat.

Ainsi, le prix d'un straddle vaut environ 0.8 écart-type, et un écart-type vaut environ 1.25 fois le prix du straddle.

\iqwrong{Prendre le prix du straddle pour l'écart-type. C'est le mouvement absolu attendu, qui est plus petit.}
\end{solution}

\begin{solution}{iq:iv2:options-and-bond-arithmetic-at-speed:13}
La gamma vaut environ 0.02 par dollar. Le theta vaut environ 1.6 cent par jour. Pour le compenser, le titre doit bouger d'environ 1.26~\$, soit environ 1.26~\%, ce qui correspond au mouvement quotidien à une volatilité de 20~\%.

La gamma\index{gamma} vaut $\varphi(d_1)/(S\sigma\sqrt T) \approx 0.4/(100 \times 0.2) = 0.02$. La valeur exacte est d'environ 0.0198.

La première voie vers le theta\index{theta} passe par Black-Scholes avec des taux nuls~: $\theta = -\tfrac12\Gamma S^2\sigma^2$, soit $0.5 \times 0.02 \times 10\,000 \times 0.04 = 4$ par an, ou $4/252 \approx 0.016$ par jour. La seconde voie utilise la règle $C \approx 0.4S\sigma\sqrt T$~: la dérivée en $T$ donne $C/(2T)$, soit $8/2/252 \approx 0.016$.

Le gain de la gamma pour un mouvement $m$ vaut $\tfrac12\Gamma m^2$, et en l'égalant au theta on obtient $m^2 = S^2\sigma^2/252$. Le mouvement de seuil est donc égal à l'écart-type d'une journée, $20/\sqrt{252} \approx 1.26$.

\iqwrong{Croire que le theta et la gamma sont deux faits distincts. Pour une option couverte en delta, ce sont les deux faces du même pari sur l'ampleur du mouvement quotidien.}
\end{solution}

\begin{solution}{iq:iv2:options-and-bond-arithmetic-at-speed:14}
La différence entre le call et le put vaut environ 4.90~\$. Sans le dividende, elle vaudrait environ 6.88~\$.

Avec un dividende, $C - P = S - De^{-rt} - Ke^{-rT}$. Le dividende est actualisé sur un trimestre à 4~\%, soit un facteur d'environ 0.99, donc environ 1.98. Le strike est actualisé sur un semestre, soit un facteur d'environ 0.98, donc environ 93.12. Alors $100 - 1.98 - 93.12 = 4.90$. Sans le dividende, on obtient $100 - 93.12 = 6.88$.

\iqwrong{Retrancher le dividende en totalité, sans l'actualiser. L'erreur est petite ici, mais la plus grave consiste à l'oublier lorsque les options franchissent la date~: la parité est alors fausse de près de 2~\$.}
\end{solution}

\begin{solution}{iq:iv2:options-and-bond-arithmetic-at-speed:15}
Les options impliquent un forward d'environ 100.51~\$. Le seul portage implique 101.50~\$. Il existe donc un écart de 1~\$, soit environ 2~\% par an, qui peut provenir d'un dividende, d'un coût d'emprunt du titre ou d'une erreur sur le taux.

Acheter un call et vendre un put de strike $K$ revient à acheter à terme au prix $K$, donc le forward moins $K$, actualisé, vaut $C - P$. Avec des taux simples, $0.50 \times 1.015 \approx 0.51$, si bien que le forward vaut environ 100.51~\$. Le seul portage donne $100 \times (1 + 3~\% \times 0.5) = 101.50$.

Le titre se négocie à un forward plus bas que ce que prévoit le portage. Quelqu'un accepte donc de payer moins cher pour détenir le titre à terme que ce que coûte son portage. Si l'on peut emprunter le titre à bas coût et qu'il n'y a pas de dividende, on peut acheter la synthèse et vendre le titre.

\iqwrong{Parler immédiatement d'arbitrage. Un écart de cette taille correspond en général à un dividende oublié ou à un titre difficile à emprunter, et sur un titre difficile à emprunter, l'arbitrage ne peut pas se faire à bon compte.}
\end{solution}

\begin{solution}{iq:iv2:options-and-bond-arithmetic-at-speed:16}
Son prix vaut 100 parce que le coupon est égal au rendement. Sa duration modifiée est égale au facteur d'annuité $(1 - 1.05^{-10})/0.05 \approx 7.72$. Le prix baisse donc d'environ 0.77 pour atteindre environ 99.23, et le prix exact vaut 99.23.

Un coupon versé au même taux que celui d'actualisation rapporte exactement ce que rapporte le principal, si bien que l'obligation vaut le pair. En dérivant le prix par rapport au rendement, on obtient, au pair, la valeur actuelle d'une annuité de 1. La duration modifiée d'une obligation au pair est donc le facteur d'annuité.

\iqwrong{Prendre 10 comme duration. Les flux d'un obligataire arrivent avant l'échéance, si bien que sa duration est plus courte, et elle diminue lorsque le coupon augmente.}
\end{solution}

\begin{solution}{iq:iv2:options-and-bond-arithmetic-at-speed:17}
Le portage vaut 0.25 point et le roll-down environ 0.11 point, soit un total d'environ 0.36. Le rendement peut monter d'environ 8 points de base avant que la position ne perde de l'argent.

Le carry\index{carry} correspond au rendement moins le taux de financement, multiplié par le temps~: $(4.5\% - 3.5\%) \times 0.25 = 0.25$. Le roll-down\index{roll-down} provient du vieillissement. Au bout de trois mois, l'obligation a 4.75 ans à courir et, sur une courbe inchangée, elle rapporte 2.5 points de base de moins. Cela rapporte $4.5 \times 0.025 = 0.1125$.

La hausse de seuil vaut $0.3625/4.5 \approx 8$ points de base.

\iqwrong{Compter le coupon comme portage sans retrancher le financement. C'est le financement qui fait qu'une position sans coût de capital ne vaut rien.}
\end{solution}

\begin{solution}{iq:iv2:options-and-bond-arithmetic-at-speed:18}
La valeur publiée fait apparaître une prime de 1.25~\%. La prime réelle est d'environ 0.25~\%, sur une juste valeur de 40.40~\$. Votre coût est de 20 points de base, il reste donc environ 5 points de base, ce qui constitue au mieux une opération marginale.

Tokyo est fermée, si bien que la valeur liquidative publiée\index{valeur liquidative} est périmée. Les contrats à terme ont bougé de 1~\%, donc la meilleure estimation du panier vaut $40 \times 1.01 = 40.40$. Alors $40.50/40.40 - 1 \approx 0.25\%$.

\iqwrong{Vendre le ETF parce qu'il «~se négocie 1.25~\% au-dessus de sa valeur~». Le ETF intègre une information que le chiffre publié ne peut pas refléter, et il ne peut le faire que parce qu'il se négocie pendant que Tokyo est fermée.}
\end{solution}

\begin{solution}{iq:iv2:options-and-bond-arithmetic-at-speed:19}
La juste valeur est de 4\,030. Le contrat à terme est cher de 6 points, soit 15 points de base du spot. Vendre le contrat à terme et acheter le panier. Après des coûts de 10 points de base, il en reste 5, soit 2 points.

La juste valeur correspond au spot, plus le financement, moins les dividendes~: $4000 \times (1 + (5\% - 2\%) \times \tfrac14) = 4030$. L'écart de 6 points vaut $6/4000 = 15$ points de base. On porte ce cash and carry\index{cash and carry} jusqu'à l'échéance, et le contrat à terme et le panier convergent alors.

\iqwrong{Donner l'écart sans ses coûts ni ses risques. Les dividendes sont des estimations, et un décalage de ceux-ci ou de votre financement peut être plus grand que les 5 points de base qui restent.}
\end{solution}

\begin{solution}{iq:iv2:options-and-bond-arithmetic-at-speed:20}
Le prix indique une probabilité d'environ 37~\% de finir au-dessus d'environ 105~\$. La valeur de Black-Scholes est d'environ 36.5~\%.

Un call spread rapporte 0 en dessous de 100 et la pleine valeur de 10 au-dessus de 110, avec une variation linéaire entre les deux. Sa valeur par unité de largeur est donc à peu près la probabilité risque-neutre de finir au-delà du milieu, c'est-à-dire d'une option digitale\index{option digitale} en 105. Et $3.67/10 = 0.367$. Pour la valeur du modèle, $d_2 = \ln(100/105)/0.2 - 0.1 = -0.344$, et $\Phi(-0.344) \approx 36.5\%$.

On lit donc une probabilité à partir de deux prix de call.

\iqwrong{Y voir la probabilité réelle. Elle est risque-neutre, et elle intègre les primes de risque ainsi que le smile de volatilité.}
\end{solution}

\begin{solution}{iq:iv2:options-and-bond-arithmetic-at-speed:21}
L'estimation normale donne environ 3.96~\$, et Black-Scholes environ 4.29~\$. Elle est plus basse parce qu'une action lognormale a une queue de droite plus épaisse qu'une loi normale de même écart-type.

La formule normale, le modèle de Bachelier\index{modèle de Bachelier}, s'écrit $C = s\,(\varphi(d) - d(1 - \Phi(d)))$, avec $d = (K - S)/s = 0.5$ et $s = 20$. On obtient alors $20 \times (0.35 - 0.5 \times 0.31) \approx 3.9$, et avec quatre chiffres significatifs, 3.96.

La valeur exacte de Black-Scholes est 4.29. L'écart d'environ 8~\% montre à quel point une option d'un an à 20~\% s'éloigne de la limite normale.

\iqwrong{Tracer une droite à partir du prix à la monnaie en utilisant son delta d'environ un demi~: $8 - 5 = 3$. Un call est convexe en fonction du strike, si bien que sa valeur baisse de plus en plus lentement à mesure qu'on s'éloigne de la monnaie, et la droite donne un résultat trop bas.}
\end{solution}

\begin{solution}{iq:iv2:options-and-bond-arithmetic-at-speed:22}
La règle empirique donne environ 0.16, puisque $\Phi(-1) \approx 0.16$. Le delta exact vaut environ 0.21 avec un an et 20~\%, et environ 0.17 pour l'option courte à 5~\%. Ils diffèrent parce qu'une loi lognormale a son centre et sa largeur sur une échelle différente de celle d'une loi normale, et cet écart compte davantage à mesure que la volatilité totale augmente.

Le delta vaut $\Phi(d_1)$, avec $d_1 = \ln(S/K)/(\sigma\sqrt T) + \tfrac12\sigma\sqrt T$. À 120 et 20~\%, $d_1 = \ln(100/120)/0.2 + 0.1 \approx -0.81$, ce qui donne 0.21. À 105 et 5~\%, $d_1 = \ln(100/105)/0.05 + 0.025 \approx -0.95$, ce qui donne 0.17.

Le logarithme de 1.2 ne vaut que 0.182, et $\tfrac12\sigma\sqrt T$ l'augmente encore. Ainsi, pour une volatilité totale plus élevée, un strike situé à un écart-type en dollars est plus proche qu'il ne l'est en logarithmes. La règle convient aux faibles volatilités totales.

\iqwrong{Utiliser la règle pour fixer un strike en vue d'une couverture précise. C'est un repère grossier pour les échéances courtes.}
\end{solution}

\begin{solution}{iq:iv2:options-and-bond-arithmetic-at-speed:23}
Vendez environ 12.1 millions de dollars du 2 ans et environ 2.77 millions de dollars du 10 ans. Le spread du butterfly est de 20 points de base. Si seul le 5 ans baisse de 5 points de base, vous gagnez 23\,000~\$. Si les trois taux montent de 10 points de base, vous ne gagnez rien, au premier ordre.

La position sur le 5 ans a un DV01 de $10 \times 460 = \$4\,600$. Chaque aile en reçoit la moitié~: $2\,300/190 \approx 12.1$ millions, et $2\,300/830 \approx 2.77$ millions. Le net est nul, si bien que les mouvements parallèles se compensent. Répartir la moitié sur chaque aile signifie aussi qu'un pur pentage se compense, lorsque les ailes bougent en sens opposé.

Le spread vaut $2 \times 4.4 - 4.0 - 4.6 = 0.2$. Si le 5 ans seul baisse, vous gagnez $4\,600 \times 5$.

\iqwrong{Égaliser les notionnels au lieu des risques. Vendre 5 millions de dollars de chaque aile laisse $4\,600 - 950 - 4\,150 = -500$, soit un DV01 net vendeur de 500~\$, ce qui constitue un pari sur le niveau de la courbe.}
\iqfollow{Cette opération est-elle neutre en trésorerie~?}{Non. Vous achetez 10 millions de dollars et vendez environ 14.9 millions de dollars de valeur nominale, si bien qu'il vous faut financer la différence.}
\end{solution}

\begin{solution}{iq:iv2:options-and-bond-arithmetic-at-speed:24}
Vous encaissez environ 6.91~\$. Un mouvement de 1~\% par jour correspond à une volatilité d'environ 16~\%. Votre profit espéré est d'environ 3.2~\$. C'est le straddle à 30~\% moins le straddle à 16~\%, et non un chiffre tiré d'une gamma constante.

Le straddle vaut $0.8 \times 100 \times 0.3 \times \sqrt{1/12} \approx 6.93$, et la valeur exacte est 6.91. À 16~\%, le straddle vaut environ 3.68. Leur différence est d'environ 3.2.

Une règle indique que la vega multipliée par l'écart de volatilité donne le résultat. La vega du straddle vaut $0.8 \times 100 \times 0.2887 = 23.1$ par 1.00, et multipliée par 0.14 elle donne 3.23. Cela tient parce qu'à la monnaie, le prix est presque linéaire en volatilité.

Prenons maintenant une gamma constante. La gamma du straddle vaut environ 0.092, et $\tfrac12\Gamma S^2(\sigma_i^2 - \sigma_r^2)T$ ne donne qu'environ 2.47. Cette estimation sous-estime le résultat, car la gamma n'est pas constante~: dans les derniers jours, elle augmente pour un titre qui reste proche du strike. L'énoncé exact est $V(\sigma_i) - V(\sigma_r) = E\left[\int \tfrac12\Gamma_i S^2(\sigma_i^2 - \sigma_r^2)\,dt\right]$, par le lemme d'Itô\index{lemme d'Itô} appliqué à la valeur de l'option.

Une trajectoire unique peut s'écarter fortement de cette moyenne, car le gain dépend de l'endroit où se trouve le titre lorsque les mouvements se produisent.

\iqwrong{Affirmer un profit certain de 3.2~\$. Couvrir à la volatilité implicite donne une moyenne et non un gain certain. Seule une couverture à la volatilité réalisée effective rend le profit certain, et l'on ne le sait pas au départ.}
\end{solution}

\begin{solution}{iq:iv2:options-and-bond-arithmetic-at-speed:25}
Le straddle implique un écart-type total de 10~\%. Cinq jours à 30~\% ont une variance d'environ $0.0018$, ce qui laisse un mouvement de résultats d'environ 9.1~\%. Si vous pensez que le mouvement sera d'environ 6~\%, vous vendez le straddle et vous couvrez.

L'écart-type total vaut $8/0.8 = 10$ points, soit une variance de 0.01. Chaque jour ordinaire a une variance de $0.3^2/252 \approx 0.000357$, soit 0.00179 pour cinq jours. Le reste vaut $0.00821$, et sa racine est d'environ 9.06~\%.

La variance\index{volatilité d'événement} s'additionne. Les écarts-types ne s'additionnent pas, si bien que retrancher les 4.2~\% de cinq jours à 10~\% donnerait 5.8~\%, ce qui est faux.

À votre hypothèse de 6~\%, la variance totale vaut $0.0036 + 0.00179 = 0.00539$, soit un écart-type de 7.34~\%. Un straddle juste vaut alors environ $0.8 \times 7.34 \approx \$5.87$, contre les 8~\$ auxquels vous pouvez vendre.

\iqwrong{Diviser le prix du straddle par le nombre de jours, ou oublier de retrancher les jours ordinaires. Le jour des résultats porte ici environ 82~\% de la variance.}
\iqfollow{Quels sont les risques~?}{Les mouvements historiques lors des résultats peuvent être inférieurs au prochain, et le titre peut sauter au-delà de vos couvertures. Dimensionnez la position de sorte qu'un saut de trois fois vos 6~\% reste supportable.}
\end{solution}
